
Code generation models achieve high execution correctness on benchmarks like HumanEval and MBPP, yet a critical gap persists between functional correctness and real-world code quality. Becker et al. (2025) demonstrated that AI-generated code runs 19\% slower than human-written solutions despite passing all unit tests. This efficiency gap reveals a fundamental limitation: execution-only optimization is necessary but insufficient for code quality.

Recent work has explored multi-objective approaches incorporating auxiliary metrics beyond execution correctness. CodeRL and CURE employ structural similarity proxies alongside test-driven feedback. CodeUltraFeedback and SEAlign use human preference signals. Lei Chen et al. (2025) demonstrated that multi-granularity structured reinforcement learning breaks supervised fine-tuning plateaus through dual textual and visual rewards for chart-to-code generation. However, these methods adopt auxiliary metrics without systematically validating measurement reliability---a prerequisite that existing work has overlooked.

The core problem is conceptual: reward engineering has remained heuristic art rather than validated methodology. A proxy metric that correlates with human judgment may still exhibit high measurement variance (intra-implementation instability), fail to distinguish algorithmic complexity classes (weak discriminability), or depend on hardware-specific noise (poor cross-platform generalization). Without systematic validation, researchers discover proxy failures only after expensive reinforcement learning training runs.

This work reframes proxy adoption as a measurement theory problem. We ask three fundamental questions for each candidate proxy: (1) Is the metric stable across repeated measurements of the same code? (2) Does it discriminate between complexity classes? (3) Does it generalize across hardware platforms? These criteria transform reward design from correlation-based heuristics to rigorous reliability testing.

Our proof-of-concept finding challenges common assumptions: structural metrics like CodeBLEU---deterministic functions of abstract syntax trees and dataflow graphs---exhibit measurement reliability suitable for optimization (CV=1.39\%, 72\% below the 5\% threshold in proof-of-concept validation). In contrast, runtime efficiency measurements using a wall-clock noise model marginally exceed the CV threshold (6.22\% vs 5.0\%), suggesting that specialized hardware instrumentation (CPU performance counters) may be required. Literature evidence from COFFE (2025) reports instruction-count CV ~2-3\% with hardware counters, though empirical confirmation in our experimental setup is needed.

This insight is compositional: different quality dimensions have distinct measurement reliability profiles. CodeBLEU's sub-metrics (n-gram match, AST match, dataflow match) are reproducible computations from static code analysis. Runtime measurements via wall-clock execution are affected by OS process scheduling, thermal throttling, and I/O latency---sources of noise that may be reduced through hardware performance monitoring units that isolate instruction count from system-dependent factors. This compositional nature means proxies must be validated independently because measurement prerequisites differ by dimension.

We make three contributions:

\textbf{1. Methodological: Four-Stage Validation Pipeline with Stage 1 Empirically Validated.} We present a framework that tests candidate proxies through increasing rigor. Stage 1 (measurement reliability: CV $\leq 5$\%, Cohen's d $\geq 0.8$, Spearman $\rho \geq 0.8)$ is empirically validated via proof-of-concept. Stages 2-4 are designed for future validation: Stage 2 (conditional independence via hierarchical regression testing $\Delta R^2 \geq 0.03)$, Stage 3 (cross-domain generalization via leave-cluster-out validation), and Stage 4 (optimization constraints via per-task execution monitoring). The pipeline employs a scoped gate design where $\geq 1$ proxy passing all Stage 1 criteria constitutes progression, preventing all-or-nothing failure. This design recognizes that partial validation---identifying which proxies are reliable---is scientifically valuable.

\textbf{2. Empirical: Validated Measurement Profiles for Code Quality Dimensions.} Proof-of-concept validation demonstrates that structural similarity (CodeBLEU) achieves exceptional measurement reliability (CV=1.39\%, Cohen's d=4.51, Spearman $\rho$=0.949), passing all three Stage 1 criteria with substantial margins. The low coefficient of variation reflects CodeBLEU's deterministic computational basis: AST parsing and dataflow analysis yield identical results when applied to the same code. The large effect size demonstrates that CodeBLEU discriminates algorithmic complexity classes. In contrast, runtime efficiency measurements using a wall-clock noise model marginally fail the CV threshold (6.22\% vs 5.0\%), despite passing Cohen's d (1.77) and Spearman $\rho$ (0.999) criteria. Literature evidence suggests CPU instruction counting via hardware performance counters should achieve lower CV, though empirical confirmation is needed.

\textbf{3. Practical: Scoped Gate Design Enabling Incremental Progress.} The validation framework employs a MUST\_WORK gate with scoped success criteria: if $\geq 1$ proxy passes all Stage 1 tests, the hypothesis proceeds with a reduced but validated proxy set. When proof-of-concept validation identified CodeBLEU as reliable while flagging efficiency measurement prerequisites, the gate returned ''PARTIAL PASS'' with actionable recommendations: continue downstream analyses with CodeBLEU alone while re-implementing runtime measurements with hardware instrumentation in parallel.

These contributions converge on a central thesis: before optimizing for multi-dimensional code quality, we must validate that quality dimensions are measurable. This work establishes construct validation as a prerequisite for proxy-based optimization.
